% !TeX root = xpsejal00-gcc-plugin.tex
\documentclass[xpsejal00-gcc-plugin.tex]{subfiles}
\begin{document}

\chapter{Implementation}\label{chap:implementation}

Building on the conceptual and architectural design outlined previously, this chapter details the practical implementation of the modernized Code Listener.
It emphasizes the engineering decisions, programming techniques, and build systems that translated the object-oriented design into a functional prototype.

The structure of this chapter follows the software's logical execution path.
Section~\ref{sec:dev-env} addresses the development environment and build artifacts, while Section~\ref{sec:build-system} describes the physical source code organization.
Subsequent sections cover the live GCC plugin pipeline (Section~\ref{sec:gcc-plugin-pipeline}), the compiler-independent core and analyzer interfaces (Section~\ref{sec:core-representation}), and the offline serialization path (Section~\ref{sec:serialization}).
The chapter concludes with implemented testing and tooling support (Section~\ref{sec:testing}), current extension points and limitations (Section~\ref{sec:implementation-limitations}), and the size of the authored implementation (Section~\ref{sec:implementation-size}).
In accordance with software engineering best practices, exhaustive source code listings are omitted from the main text and are instead provided on the attached digital medium or on GitHub~\footnote{\url{https://github.com/pseja/gcc-plugin-for-static-analyzer-support}}.

\section{Development Environment and Platform Constraints}\label{sec:dev-env}

Developing a compiler plugin imposes strict platform and toolchain constraints.
Like the original Predator toolchain, the modernized infrastructure targets Linux environments.
Because the plugin is compiled as a shared object and loaded directly into the GCC process memory space, it is intrinsically tied to the GCC Plugin API.\@
The build system must match the exact GCC installation used at runtime and compile against that specific version's plugin headers.
This constraint is managed via a \texttt{TARGET\_GCC} CMake variable.
By default, it is set to \texttt{gcc-12}.
If the required compiler headers are missing during configuration, the build is designed to abort immediately rather than producing a partially or silently broken binary.

The implementation utilizes C++23, the latest stable version of the C++ standard available during development~\cite{gccCxxStatus}.
This choice enables the internal models to leverage standard library facilities, thereby avoiding reliance on complex, custom framework code.
For instance, \texttt{std::variant} was employed for closed type and instruction hierarchies, \texttt{std::optional} was used to safely handle incomplete compiler metadata, and modern range-based traversal was applied where appropriate within the Code Model.

To ensure safety and compatibility, the generated binaries are compiled conservatively.
All plugin-facing targets are compiled with \texttt{-fPIC} to support dynamic loading.
Furthermore, the codebase uses \texttt{-fno-rtti} to avoid the overhead of runtime type information, and enforces strict compilation checks via \texttt{-Wall}, \texttt{-Wextra}, \texttt{-Werror}, and \texttt{-pedantic} to catch potential defects at compile time.

External dependencies were kept to an absolute minimum.
Persistent JSON storage is handled by the header-only \texttt{nlohmann/json} library, avoiding the introduction of shared-library dependencies that could cause symbol collisions inside the GCC process~\cite{nlohmannJson}.
Additionally, the standard \texttt{patch} utility is required during the build configuration to automatically apply compatibility shims to the legacy Predator codebase, ensuring the experiment is reproducible from a clean repository checkout.

\section{Build System and Source Code Organization}\label{sec:build-system}

A two-layered build system is employed in the implementation.
CMake manages the complex target graph, dependencies, and test registrations, while a top-level Makefile offers concise commands for standard development workflows, including \texttt{make build}, \texttt{make test}, \texttt{make json}, \texttt{make callgraph}, and others.

The build process produces five distinct categories of artifacts:

\begin{enumerate}
  \item \texttt{cl\_core}: The compiler-independent static library housing the Code Model.
  \item \texttt{cl\_gcc}: The shared library that GCC dynamically loads as the actual plugin.
  \item \textbf{Analyzer Libraries:} Dedicated shared libraries including \texttt{cl\_callgraph\_dot} and \texttt{cl\_recursion\_check}, with additional analyzers built in the same way.
  \item \textbf{Standalone Executables:} The offline tools \texttt{cl\_merge} and \texttt{cl\_analyze}.
  \item \textbf{Auxiliary Outputs:} Generated Doxygen documentation, class diagrams, and CTest registrations.
\end{enumerate}

To maintain the architectural boundaries defined in the design phase, the physical directory layout strictly mirrors the separation of concerns:

\begin{itemize}
  \item \texttt{src/core/}: Contains the compiler-independent Code Model, entity pools, and analyzer interfaces.
  \item \texttt{src/adapters/gcc/}: Houses the GCC plugin registration, the GIMPLE translation logic, and argument parsing.
  \item \texttt{src/adapters/predator/}: Contains the legacy compatibility layer for Predator.
  \item \texttt{src/annotation\_services/}: Implements the lazily evaluated analyses, most importantly the Call Graph service used by multiple analyzers.
  \item \texttt{src/exporters/}: Contains the DOT, JSON, and pretty-printed export pipelines, together with the JSON frontend and merger.
  \item \texttt{src/tools/}: Holds the entry points for the \texttt{cl\_merge} and \texttt{cl\_analyze} executables.
  \item \texttt{include/}: Exposes the C and C++ analyzer Application Binary Interfaces (ABIs).
  \item \texttt{tests/}: Contains the native regression test suites.
\end{itemize}

This directory structure enforces the most critical implementation rule of the entire project: GCC headers must never leak into compiler-independent code.
All direct dependencies on the GCC API are quarantined entirely within \texttt{src/adapters/gcc/}.
This physical boundary makes it possible to compile the offline \texttt{cl\_analyze} tool and the Code Model without requiring the GCC runtime data structures.

\section{Realizing the GCC Plugin Pipeline}\label{sec:gcc-plugin-pipeline}

The live compiler plugin represents the first of the three primary execution paths.
Its entry point was implemented in \texttt{register\_plugin.cpp}.
Its sole responsibility is to verify that the runtime GCC version matches the compiled version before delegating initialization to the \texttt{PluginContext} singleton.
Inside \texttt{PluginContext.cpp}, the plugin registers its custom GIMPLE pass immediately following GCC's native \texttt{cfg} pass.
It records the current input file at the \texttt{PLUGIN\_START\_UNIT} event and finally installs a \texttt{PLUGIN\_FINISH} callback that triggers the model export and analyzer invocation at the end of the compilation unit.

Command-line argument parsing is centralized in \texttt{PluginArgs.cpp}.
It parses the raw GCC plugin arguments, converts typed values, like exporter verbosities, and validates them.
The arguments currently wired into runtime behavior cover verbosity selection, analyzer loading, forwarded analyzer arguments, and the JSON, DOT, and PP export requests.
The parser also recognizes flags intended as future work hooks, such as \texttt{help}, \texttt{version}, \texttt{dry-run}, \texttt{dump-types}, \texttt{preserve-ec}, \texttt{pid-file}, and \texttt{type-dot}, preventing the need to hardwire ad hoc conditionals throughout the plugin logic later.

The custom GCC pass is defined in \texttt{Pass.cpp}.
For every function, its execute method forwards work directly to \texttt{GCCAdapter::processFunction}.
\texttt{GCCAdapter.cpp} is the largest single implementation file in the project.
It incrementally populates the new Code Model by traversing the GCC trees and GIMPLE statements.
To optimize this translation:

\begin{itemize}
  \item Compiler types are memoized using a cache keyed by the GCC tree nodes, ensuring that recursive types are represented only once.
  \item Anonymous compiler aggregates are assigned synthetic location-based names, significantly improving the readability of diagnostic and JSON outputs.
  \item Source locations are aggressively normalized via helper functions to recover coordinates even for GIMPLE statements missing complete metadata, providing the best possible information for diagnostics and analysis.
  \item The adapter stores only internal identifiers, entirely stripping away the raw GCC memory pointers so that the resulting Code Model outlives the compiler process.
\end{itemize}

Finally, diagnostic messages are routed through \texttt{GCCDiagnosticReporter.cpp}.
This class maps framework diagnostics into native GCC calls, specifically \texttt{inform}, \texttt{warning\_at}, \texttt{error\_at}, \texttt{fatal\_error}, ensuring that any bugs found by an analyzer appear to the user exactly like standard compiler errors.
Since GCC has no built-in debug-level reporting channel, the plugin provides a development-only path that prints grey \texttt{stderr} lines containing the relative file path, line number, function name, and message text.
These messages are emitted only when the verbosity level is set to 2 or higher, allowing adapter debugging to remain available on demand without affecting ordinary analyzer runs.

\section{Core Representation and Analyzer Interfaces}\label{sec:core-representation}

The internal Code Model is driven by centralized entity pools.
In the implementation, these pools for types, variables, functions, blocks, and instructions are instantiated as \texttt{std::deque} containers.
A \texttt{std::deque} allocates memory in fixed chunks, meaning that as the GCC adapter incrementally populates the model and holds temporary references, the memory addresses of previously inserted entities remain perfectly stable~\cite{boostedCppDeque}.
This avoids the pointer-invalidation issues common to \texttt{std::vector} reallocations.

Because deep inheritance trees were avoided, instruction and type payloads are implemented using C++23 \texttt{std::variant}.
This allows each instruction to be stored as a plain value type.
Later processing, such as exporting or analyzing the model, is performed cleanly by branching over the closed set of alternatives using \texttt{std::visit}.

To load the analysis logic, \texttt{PluginContext::load\_analyzer} performs dynamic library loading.
It explicitly checks the library for the modern \texttt{cl\_get\_native\_api} symbol.
If the symbol exists and versions match, the library is wrapped in the \texttt{NativeAnalyzerBridge}.
If the modern API is missing, the infrastructure falls back to the historical C ABI through the legacy \texttt{cl\_get\_analyzer\_api} entry point.
The resulting listener is then wrapped in a \texttt{LegacyPredatorBridge}.
Missing symbols or version mismatches are reported as errors via the plugin diagnostic channel and prevent the analyzer from loading.

Modern native analyzers are provided with an \texttt{AnalysisContext} declared in the public analyzer headers.
Its implementation is located in \texttt{src/exporters/AnalysisContext.cpp}.
This context provides three specific services: the diagnostic reporter, the \texttt{AnalysisManager}, and unified export helpers exposed through \texttt{exportDot}, \texttt{exportPP}, and \texttt{exportJson}.
Because the \texttt{AnalysisManager} uses lazy evaluation, multiple analyzers that run sequentially can request and share expensive derived data, such as the Call Graph, from a single cached instance.

Conversely, backward compatibility is handled by \texttt{PredatorAdapter.cpp}.
This adapter systematically walks the modern Code Model and reconstructs the legacy \texttt{cl\_type} and \texttt{cl\_var} structures.
The emission sequence was carefully implemented to match historical expectations: all types and variables are pre-registered before any function bodies are streamed.
This exact ordering ensures the legacy Predator verification kernel functions entirely unmodified.

\section{Serialization and Offline Whole-Program Workflow}\label{sec:serialization}

The second major execution path is the offline workflow, which completely decouples the analysis from the live compiler process.
Serialization is implemented via the \texttt{JSONExporter} and \texttt{JSONImporter} classes.
Because the model relies solely on strongly typed IDs rather than memory addresses, the serialization process maps internal IDs to JSON values almost directly, avoiding complex graph-pointer reconstruction.

Whole-program analysis requires merging multiple translation units into a single model.
This is executed by \texttt{ModelMerger.cpp}.
To prevent ID collisions, for example, \texttt{FunctionId(0)} existing in two separate files, the merger generates explicit remapping tables.
As a new file is appended to the consolidated Code Model, the merger systematically rewrites every cross-reference through these tables, seamlessly preserving the integrity of the graph topology across file boundaries.

These capabilities are exposed to the user via two standalone executables.
The \texttt{cl\_merge} tool accepts multiple \texttt{.json} files, leverages the \texttt{ModelMerger}, and outputs a unified program model.
The \texttt{cl\_analyze} tool acts as the offline analysis driver.
It can either load a native analyzer library or invoke the built-in offline exporters for DOT and pretty-printed output.
Instead of using the \texttt{GCCDiagnosticReporter}, it instantiates an equivalent \texttt{StderrDiagnosticReporter}.
When the analyzer path is selected, it loads the target analyzer and feeds it the deserialized JSON model.
Consequently, the analyzer backend requires no modifications to operate offline; the semantic contract remains identical whether the model originated from GCC memory or a JSON text file.

\section{Testing, Tooling, and Practical Workflow}\label{sec:testing}

To ensure the architecture serves as a robust research prototype, extensive automated testing and tooling were implemented.
This section describes the implemented testing infrastructure and practical usage flow, while the quantitative evaluation itself is deferred to the following testing chapter.
When the \texttt{WITH\_PREDATOR} CMake option is enabled, the build system compiles the patched Predator components and exposes their regression suite through CTest.

To verify correctness across the different architectural boundaries, CTest is configured with four separate regression suites:
\begin{enumerate}
  \item \textbf{GCC-Adapter Tests}: 46 tests validating that the GIMPLE translation produces the correct JSON output against historical fixtures.
  \item \textbf{Code Model Integration Tests}: 38 tests focusing specifically on types, initializers, control flow edges, and edge cases.
  \item \textbf{Offline Pipeline Tests}: 3 end-to-end tests covering the offline analysis path through the \texttt{cl\_analyze} tool.
  \item \textbf{Predator Regression Suite}: 1,192 tests executed against the patched Predator kernel to validate backward compatibility on the preserved listener interface.
\end{enumerate}

To demonstrate how the system is used in practice, two reference native analyzers were implemented: \texttt{callgraph\_dot} and \texttt{recursion\_check}.
The \texttt{callgraph\_dot} analyzer requests the shared \texttt{CallGraph} annotation and emits its nodes and edges as a Graphviz DOT file.
The \texttt{recursion\_check} analyzer requests the same annotation, iterates over its strongly connected components, and reports direct and mutual recursion both to diagnostics and to a text report.
A developer can run \texttt{make callgraph FILE=example.c} to generate a \texttt{callgraph.dot} visualization of the program's interprocedural structure.
Following that, \texttt{make recursion FILE=example.c} checks the same program for direct and mutual recursion, and writes \texttt{recursion\_report.txt} by default.
In both cases, the output path can be redirected via the forwarded-analyzer argument string exposed by the top-level Makefile.
Together, these analyzers show that both simple visual tools and more analysis-oriented native analyzers can be written over the new architecture with minimal boilerplate.

At this stage, the implementation already delivers three concrete execution paths: the live GCC plugin path, the Predator compatibility path, and the offline JSON-based path.
The subsequent testing chapter quantitatively evaluates these execution paths, whereas the present chapter remains focused on their construction.

\section{Current Limitations and Extension Points}\label{sec:implementation-limitations}

While the architecture is stable, several limitations and extension points remain.
The \texttt{AnalysisContext} and \texttt{AnalysisManager} are explicitly designed as extension points for future annotation services.
However, the offline \texttt{cl\_analyze} workflow currently only supports native analyzers or built-in exporters.
Executing legacy analyzers like Predator offline is not supported because they depend heavily on the \texttt{PredatorAdapter} reconstructing the legacy data stream, a process currently tied directly to the live GCC execution path.

These limitations do not diminish the prototype's practical value; rather, they delineate the next logical steps in its development.
These include expanding the set of shared annotations, implementing concrete behavior for currently reserved plugin arguments, and extending the offline path where feasible without compromising the existing analyzer boundary.

\section{Implementation Size}\label{sec:implementation-size}

To quantify the scale of the implementation and distinguish original thesis work from legacy code, the authored C and C++ implementation files comprise approximately 11,309 lines of code (LOC)~\footnote{Measured in April 2026}.
The repository includes an additional 2,061 lines of automated test code, bringing the total authored contribution to 13,370 lines.

\begin{table}[H]
  \begin{tabular}{ll}
    Subsystem                                                & LOC    \\
    Core model (\texttt{src/core/})                          & 2,082  \\
    GCC integration (\texttt{src/adapters/gcc/})             & 2,759  \\
    Predator compatibility (\texttt{src/adapters/predator/}) & 1,239  \\
    Annotation services (\texttt{src/annotation\_services/}) & 599    \\
    Exporters and offline frontend (\texttt{src/exporters/}) & 3,825  \\
    Command-line tools (\texttt{src/tools/})                 & 368    \\
    Native analyzers (\texttt{analyzers/})                   & 279    \\
    Public headers (\texttt{include/})                       & 158    \\
    Total implementation code                                & 11,309
  \end{tabular}
  \caption{Lines of code (LOC) for each major subsystem of the implementation.}\label{tab:loc}
\end{table}

In contrast, the imported Predator submodule comprises approximately 673,000 lines of legacy C and C++ code, with Code Listener itself totaling about 25,000 lines, including headers, support files, and tests.
This legacy code is treated solely as an external dependency and evaluation baseline, and is excluded from the thesis implementation metrics.
As illustrated in Table~\ref{tab:loc} on page~\pageref{tab:loc}, the code-size distribution reflects the architectural priorities of the project, with the largest components allocated to the Exporter pipelines, the GCC translation adapter, and the Core representation.

\end{document}
